\section{Paso 1: formulación del problema y planificación}
\label{sec:paso1-formulacion}
% Origen: 01-definicion-del-caso.md §1, §3, §5 y §6; 07-plan-tpi.md §3 (D1, D4, D5).

\subsection{Problema y pregunta de decisión}

En el pico de la mañana, el vestíbulo Principal de Constitución recibe oleadas de pasajeros del Roca. Por
ventana se observan en uso unos 19,5 molinetes de los 22 que el dataset de SBASE identifica en ese
vestíbulo. El resto de la capacidad instalada no se usa o se usa muy poco. Además, la carga entre
molinetes es muy desigual: el más usado procesa 11 veces más pasajeros que el menos usado
\parencite{sbase2026molinetes}. La pregunta de decisión es:

\begin{quote}
¿Cuánto se reduce la espera de los pasajeros si, en los días hábiles de 07:00 a 09:30, se habilitan los
22 molinetes del vestíbulo Principal en lugar de los 20 que hoy operan en promedio?
\end{quote}

La propuesta original hablaba de habilitar \enquote{los 28 molinetes instalados}. Al procesar el dataset
por vestíbulo se comprobó que 28 es el total de la estación: 22 molinetes en el Principal y 8 en el
vestíbulo Plaza (sección~\ref{sec:paso2-datos-modelo-conceptual}). La pregunta se corrigió en
consecuencia. La ampliación a 28 molinetes en el Principal se conserva solo como variante hipotética.

\subsection{Objetivos}

\textbf{Objetivo general.} Cuantificar, con un modelo de simulación peatonal, el efecto de la cantidad de
molinetes habilitados sobre la espera de los pasajeros que ingresan a la Línea~C por el vestíbulo
Principal de Constitución durante el pico de la mañana.

\textbf{Objetivos específicos.}
\begin{enumerate}
  \item Caracterizar la demanda por ventana de 15 minutos con el dataset abierto de SBASE y la estructura
        de las oleadas con el horario oficial del Roca.
  \item Construir y verificar en AnyLogic un modelo que represente la caminata, la elección de molinete,
        la cola y la validación de cada pasajero.
  \item Comparar el escenario actual con el alternativo mediante réplicas apareadas e intervalos de
        confianza para la diferencia de medias.
  \item Explicitar qué parámetros son dato y cuáles son supuestos, y en qué medida las conclusiones
        dependen de estos últimos.
\end{enumerate}

\subsection{Alcance}

\begin{itemize}
  \item \textbf{Sistema.} El vestíbulo Principal de la estación Constitución de la Línea~C. Para cada
        pasajero, el sistema empieza cuando entra al vestíbulo, desde el hall del Roca o desde la calle.
        Termina cuando valida y cruza a la zona paga, hacia las escaleras a los andenes.
  \item \textbf{Período.} Días hábiles, de 07:00 a 09:30. El pico está entre las 08:15 y las 08:45. La
        tarde no se modela: en Constitución el flujo vespertino es de salida y los molinetes de ingreso
        apenas llegan a 500 pasajeros por ventana.
  \item \textbf{Fuera del alcance.} El vestíbulo Plaza, los andenes y el viaje en subte, la venta y carga
        de tarjetas, los egresos y los días no hábiles.
\end{itemize}

\subsection{Medidas de rendimiento}

Sea $W_i$ la espera del pasajero $i$: el tiempo desde que se incorpora a la cola de un molinete hasta que
empieza a validar. Quien encuentra el molinete libre tiene $W_i = 0$. La caminata por el vestíbulo no
cuenta como espera. Sean $n$ los pasajeros que ingresan antes de las 09:30 y $Q(t)$ la cantidad de
personas en cola en el instante $t$, con $t$ en segundos desde las 07:00. Todas las medidas se calculan
dentro de cada réplica (tabla~\ref{tab:medidas}).

\begin{table}[H]
\centering
\caption{Medidas de rendimiento por réplica}
\label{tab:medidas}
\small
\begin{tabularx}{\textwidth}{@{}l X l@{}}
\toprule
Medida & Definición & Tipo \\
\midrule
P90 de espera & Percentil 90 de $\{W_i\}$ por rango más próximo: el valor de orden $\lceil 0{,}9\,n \rceil$ & Primaria \\
Proporción con espera $> 30$~s & $\frac{1}{n}\sum_i \mathbf{1}\{W_i > 30\}$ & Primaria \\
P90 de la cohorte pico & P90 de los pasajeros que ingresan entre $t = 4500$ y $t = 6300$ (08:15 a 08:45) & Primaria \\
Espera media & $\bar W = \frac{1}{n}\sum_i W_i$, en la franja y en la cohorte pico & Secundaria \\
Largo medio de cola & $L_q = \frac{1}{9000}\int_0^{9000} Q(t)\,dt$ & Secundaria \\
Cola máxima & $\max_{t < 9000} Q(t)$ & Secundaria \\
Utilización & Segundos ocupados antes de las 09:30 sumados sobre los molinetes, divididos por $9000 \cdot c$ & Secundaria \\
Dispersión de utilización & Desvío poblacional de las utilizaciones de los $c$ molinetes habilitados & Secundaria \\
Tiempo de drenaje final & Desde el último ingreso hasta que sale el último pasajero & Secundaria \\
Conservación & Generados $=$ procesados al terminar el drenaje & Control \\
\bottomrule
\end{tabularx}
\end{table}

El P90 se eligió además de la media porque el fenómeno es de picos: una media baja puede esconder colas
largas para los pasajeros que llegan con cada tren. La proporción sobre 30~s traduce la espera a un nivel
de servicio fácil de comunicar. La cohorte pico aísla la media hora más exigente. La cantidad de
procesados se usa como control y no como primaria: por conservación, es igual en los dos escenarios de
cada par y su diferencia vale cero por construcción.

\subsection{Escenarios}

La tabla~\ref{tab:escenarios} resume los escenarios definidos. La comparación obligatoria del estudio es
\EO{} contra \EI{}.

\begin{table}[H]
\centering
\caption{Escenarios definidos y alcance de cada uno}
\label{tab:escenarios}
\small
\begin{tabularx}{\textwidth}{@{}l X X@{}}
\toprule
Escenario & Configuración & Estado en este informe \\
\midrule
\EO{} & 20 molinetes habilitados, el entero más próximo a los 19,5 con tráfico por ventana & Evaluado \\
\EI{} & Los 22 molinetes del vestíbulo Principal & Evaluado \\
E1-b & Ampliación hipotética a 28 molinetes en el mismo banco & Implementado; no se evalúa porque el espacio físico no está verificado \\
E2 & Desvío de una fracción del flujo hacia el vestíbulo Plaza & Implementado en la capa lógica; trabajo futuro (falta modelar Plaza) \\
E3 & Validación sin contacto (tarjeta EMV o QR) en lugar de SUBE & Implementado en la capa lógica; trabajo futuro (no hay tiempo de validación citable) \\
\bottomrule
\end{tabularx}
\end{table}

\EI{} es la alternativa más directa y de menor costo: no requiere obra, solo habilitar equipos que ya
existen. Que los 22 molinetes estén efectivamente disponibles para ingreso en el pico es un supuesto. Por
eso se pidió a SBASE la correspondencia entre identificadores y equipos físicos. E2 y E3 quedan fuera de
la evaluación. E2 exige un segundo circuito para Plaza, del que no hay geometría, y un supuesto sobre
cuántos pasajeros aceptarían desviarse. E3 no tiene un tiempo de validación EMV o QR publicado para la
red, y el único molinete de ese tipo está en Plaza.
\pendiente{confirmar con la cátedra la redefinición de \EI{} (22 en lugar de 28) y el alcance de E2 y E3; consulta en 08-consulta-catedra.md.}

\subsection{Plan de trabajo}

El estudio se organizó en fases, cada una con un entregable verificable (tabla~\ref{tab:plan}).

\begin{table}[H]
\centering
\caption{Plan de trabajo}
\label{tab:plan}
\small
\begin{tabularx}{\textwidth}{@{}l X l@{}}
\toprule
Fase & Contenido & Fecha \\
\midrule
Definición & Aprobación del tema, perfilado del dataset de SBASE y definición del caso & 20 a 23/09/2026 \\
Pedidos de datos & Correos a SBASE, Emova y Trenes Argentinos; pedidos formales por Ley~104 y Ley~\num{27275} & 23/09 y 05/10/2026 \\
Modelo & Capa lógica, capa peatonal, demanda por horario del Roca y automatización de las corridas & 20/09 a 24/09/2026 \\
Pilotos & Diagnóstico del atasco, corrección y calibración del escenario intermedio & 24/09 a 05/10/2026 \\
Producción & 30 pares \EO{}-\EI{} y reproducción en una segunda máquina & 05/10 y 09/10/2026 \\
Informe y video & Redacción en \LaTeX{}, presentación y grabación & \pendiente{fecha de entrega} \\
\bottomrule
\end{tabularx}
\end{table}
